\documentclass[11pt,a4paper]{article}
\usepackage[T1]{fontenc}
\usepackage[utf8]{inputenc}
\usepackage{lmodern}
\usepackage{microtype}
\usepackage[a4paper,margin=2.25cm]{geometry}
\usepackage{longtable,booktabs,array,calc}
\usepackage{fancyvrb}
\usepackage{xcolor}
\usepackage{hyperref}
\usepackage{xurl}
\usepackage{fancyhdr}
\usepackage{enumitem}
\usepackage{parskip}
\usepackage{titlesec}
\definecolor{ddtadark}{RGB}{30,38,48}
\definecolor{ddtablue}{RGB}{38,82,126}
\hypersetup{colorlinks=true,linkcolor=ddtablue,urlcolor=ddtablue,citecolor=ddtablue}
\pagestyle{fancy}
\fancyhf{}
\lhead{Documentation-Driven Threat Analysis}
\rhead{BA4-T2 projection pressure test}
\cfoot{\thepage}
\setlength{\headheight}{14pt}
\setlist[itemize]{leftmargin=1.5em,itemsep=0.2em,topsep=0.25em}
\setlist[enumerate]{leftmargin=1.7em,itemsep=0.2em,topsep=0.25em}
\titleformat{\section}{\Large\bfseries\color{ddtadark}}{}{0em}{}
\titleformat{\subsection}{\large\bfseries\color{ddtadark}}{}{0em}{}
\titleformat{\subsubsection}{\normalsize\bfseries\color{ddtadark}}{}{0em}{}
\providecommand{\tightlist}{\setlength{\itemsep}{0pt}\setlength{\parskip}{0pt}}
\setlength{\emergencystretch}{4em}
\hbadness=10000
\hfuzz=20pt
\sloppy
\begin{document}
\section{DDTA BA4-T2 method-owned interpretation, coverage loss and
cross-projection consistency pressure test -
R1}\label{ddta-ba4-t2-method-owned-interpretation-coverage-loss-and-cross-projection-consistency-pressure-test---r1}

\textbf{Status:} COMPLETED / PROVISIONAL PASS WITH
INTERPRETATION-COVERAGE REFINEMENT

\textbf{Executed against repository baseline:}
\texttt{f90ef3a0bc0b7712cb8081165c28e8923aec9e2d}

\textbf{Input candidate:}
\texttt{BA4\_PROJECTION\_BOUNDARY\_TRACEABILITY\_CANDIDATE\_R1.md}

\subsection{1. Question under test}\label{1-question-under-test}

Can the BA4-T1 projection boundary survive two intentionally
incompatible method-oriented projections, aggressive information loss,
non-current BA state and cross-baseline rebuild without either importing
consumer semantics into Base Analysis or inventing a shared projection
ontology?

The test is falsification-first. Each probe tries to force a semantic
responsibility into the shared core. A responsibility is accepted only
when removing it creates a concrete loss of reviewability, coverage
honesty, qualification or traceability.

\subsection{2. Fixed dependencies and
guardrails}\label{2-fixed-dependencies-and-guardrails}

Fixed entering T2:

\begin{verbatim}
BA0   CLOSED
BA1   CLOSED
BA2   CLOSED
BA3   CLOSED
BA4-T1 provisional projection lower bound active
\end{verbatim}

T2 does not define:

\begin{itemize}
\tightlist
\item
  a complete STRIDE or STRIDE-AI schema;
\item
  a DFD ontology;
\item
  AnalysisRecord or Common Finding;
\item
  ThreatForge classes/tables/UI;
\item
  a universal projection taxonomy;
\item
  BA5 lexical rules; or
\item
  BA6 final holdout closure.
\end{itemize}

Method labels used below are bounded pressure instruments only.

\subsection{3. Test basis}\label{3-test-basis}

The main facial-access BA basis includes the already-reviewed meanings:

\begin{verbatim}
transfer(CameraSubsystem,
         RecognitionProcessor,
         RecognitionCapture)

correlate(RecognitionCapture, RecognitionRequest)

consumeService(project, LocalConnectivity)

assignResponsibility[negative](
  project,
  underlyingTransport,
  ownership/management)

realize(LocalConnectivity, WiredEthernet)

constrain(delivery, completion/security properties)
\end{verbatim}

Controls use facial mutations M1-M4 plus order/WMS/provider meanings,
including project-owned versus external inventory responsibility and
governed \texttt{PaymentAuthorizationResult} normalization.

\subsection{4. Probe P1 - incompatible method
taxonomies}\label{4-probe-p1---incompatible-method-taxonomies}

\subsubsection{Hypothesis attacked}\label{hypothesis-attacked}

Two method projections need a shared category ontology to remain
consistent.

\subsubsection{Projection F}\label{projection-f}

Local kinds:

\begin{verbatim}
FlowParticipant
FlowExchange
ExposureAnnotation
\end{verbatim}

It selects transfer/service-use meaning and may interpret service-use
plus negative responsibility as:

\begin{verbatim}
externally-provided-transport-exposure
\end{verbatim}

\subsubsection{Projection A}\label{projection-a}

Local kinds:

\begin{verbatim}
AssuranceSubject
AuthorityPlacement
ContractDependency
NormalizationExpectation
\end{verbatim}

It selects responsibility/contract/normalization meaning and may
interpret the same service/responsibility basis as:

\begin{verbatim}
delegated-assurance-dependency
\end{verbatim}

\subsubsection{Result}\label{result}

PASS without shared taxonomy.

The two items can be compared through overlapping role-bound BA trace.
Their category difference is method-local, not a contradiction.

\subsubsection{Disposition}\label{disposition}

\begin{verbatim}
universal shared projection ontology     REJECTED
BA trace as common semantic anchor        PASS
\end{verbatim}

\subsection{5. Probe P2 - descriptor plus trace
reproducibility}\label{5-probe-p2---descriptor-plus-trace-reproducibility}

\subsubsection{Hypothesis attacked}\label{hypothesis-attacked-1}

An immutable descriptor plus role-bound BA trace is always sufficient to
reconstruct method-owned interpretation.

\subsubsection{Counterexample}\label{counterexample}

One immutable projection revision contains two local rules whose inputs
overlap:

\begin{verbatim}
R-ext-service
  serviceUse + responsibilityBoundary
    -> externally-provided-transport-exposure

R-assurance-review
  serviceUse + responsibilityBoundary + constraintSet
    -> assurance-dependency-review
\end{verbatim}

An item with only \texttt{serviceUse} and
\texttt{responsibilityBoundary} trace identifies its source basis but
not which local rule created its method-owned meaning when several rule
paths remain valid under the same descriptor.

\subsubsection{Result}\label{result-1}

The T1 lower bound is insufficient for non-trivial method
interpretation.

\subsubsection{Smallest refinement}\label{smallest-refinement}

\begin{verbatim}
methodOwnedInterpretation present
    -> interpretationRuleRef required
\end{verbatim}

The referenced local rule revision is already contained/resolved under
the immutable projection descriptor revision.

\subsubsection{Negative control}\label{negative-control}

No BA3 derivation rule is reused. The interpretation is not
method-neutral and cannot create accepted BA meaning.

\subsection{6. Probe P3 - aggressive lossy
selection}\label{6-probe-p3---aggressive-lossy-selection}

\subsubsection{Hypothesis attacked}\label{hypothesis-attacked-2}

A generic coverage statement can represent both intentionally selective
and complete-within-scope projections.

\subsubsection{Counterexample}\label{counterexample-1}

Two views have the same eligible BA scope:

\begin{verbatim}
current accepted transfer propositions
\end{verbatim}

but one promises to render all eligible transfers while the other
intentionally chooses only transfers touching one analysis focus.

If both use only the statement
\texttt{transfer\ propositions\ are\ in\ scope}, an omitted transfer is
ambiguous:

\begin{itemize}
\tightlist
\item
  incomplete projection build; or
\item
  valid selective omission.
\end{itemize}

\subsubsection{Result}\label{result-2}

T1 coverage contract requires an explicit mode:

\begin{verbatim}
EXHAUSTIVE_FOR_DECLARED_SCOPE
SELECTIVE
\end{verbatim}

\subsubsection{Important boundary}\label{important-boundary}

An omission from an exhaustive projection is a projection completeness
defect. It is not a project-level negative proposition and does not
prove source documentation incompleteness.

\subsection{7. Probe P4 - omission as negative
fact}\label{7-probe-p4---omission-as-negative-fact}

\subsubsection{Hypothesis attacked}\label{hypothesis-attacked-3}

If an exhaustive projection omits an item, the consumer may treat the
missing meaning as false in the project.

\subsubsection{Result}\label{result-3}

REJECTED.

Even under exhaustive coverage, the claim is only:

\begin{verbatim}
this materialization failed its declared projection contract
\end{verbatim}

The source of project truth remains governed documentation through
accepted BA.

\subsection{8. Probe P5 - diagnostic/stale review
projection}\label{8-probe-p5---diagnosticstale-review-projection}

\subsubsection{Hypothesis attacked}\label{hypothesis-attacked-4}

A projection can expose stale or diagnostic BA meaning using the same
rendering semantics as current accepted BA.

\subsubsection{M4 pressure}\label{m4-pressure}

Representation context changes while an older transfer source remains
inconsistent. BA3 can mark the old transfer \texttt{STALE} and localize
a \texttt{DIAGNOSTIC\_UNRESOLVED} condition.

Two projection purposes are tested:

\begin{verbatim}
normal current-state method view
  qualificationPolicy = ACCEPTED + CURRENT

review/change view
  qualificationPolicy includes STALE / DIAGNOSTIC_UNRESOLVED
\end{verbatim}

\subsubsection{Result}\label{result-4}

The review projection is valid only when the non-current/unresolved
qualification remains visible in the consumed meaning.

The diagnostic is not promoted to project fact merely because it is
accepted as a diagnostic.

\subsubsection{Smallest refinement}\label{smallest-refinement-1}

\texttt{qualificationPolicy} is explicit inside the projection coverage
contract.

No duplicate BA3 lifecycle state machine is introduced.

\subsection{9. Probe P6 - interpretation over stale
basis}\label{9-probe-p6---interpretation-over-stale-basis}

\subsubsection{Hypothesis attacked}\label{hypothesis-attacked-5}

A method-owned annotation may be presented as a current result even when
one of its BA inputs is stale.

\subsubsection{Result}\label{result-5}

REJECTED for current shared interpretation.

A review-oriented method projection may still compute the local
annotation to identify potentially affected analysis, but it must
preserve that its basis is stale/unresolved.

No new universal projection status enum is forced; the projection
descriptor and trace to BA3-qualified inputs are sufficient.

\subsection{10. Probe P7 - cross-projection
consistency}\label{10-probe-p7---cross-projection-consistency}

\subsubsection{Hypothesis attacked}\label{hypothesis-attacked-6}

Two projections over the same BA must have equivalent categories/items
to be considered consistent.

\subsubsection{Result}\label{result-6}

REJECTED.

Consistency is checked against BA, not by pairwise method taxonomy
equality.

For the same BA baseline:

\begin{verbatim}
Projection F item
  basis = {consumeService, negative responsibility}
  interpretation = externally-provided-transport-exposure

Projection A item
  basis = {consumeService, negative responsibility}
  interpretation = delegated-assurance-dependency
\end{verbatim}

Both are valid if each is allowed by its own immutable local
interpretation rule and remains explicitly method-owned.

\subsubsection{Comparison rule}\label{comparison-rule}

\begin{verbatim}
cross-projection common denominator = BA trace
\end{verbatim}

No first-class method-to-method equivalence identity is required.

\subsection{11. Probe P8 - contradictory shared
rendering}\label{11-probe-p8---contradictory-shared-rendering}

\subsubsection{Hypothesis attacked}\label{hypothesis-attacked-7}

Different projections can disagree about shared BA truth as long as
their methods differ.

\subsubsection{Counterexample}\label{counterexample-2}

BA contains negative project responsibility for underlying transport.

Projection F renders:

\begin{verbatim}
project does not own/manage underlying transport
\end{verbatim}

Projection A renders as shared truth:

\begin{verbatim}
project owns underlying transport
\end{verbatim}

\subsubsection{Result}\label{result-7}

Projection A fails semantic preservation. Method diversity does not
excuse inversion of shared BA polarity.

No cross-projection contradiction proposition is needed; the failing
projection can be diagnosed directly against its traced BA basis.

\subsection{12. M1 replay - realization
change}\label{12-m1-replay---realization-change}

B0:

\begin{verbatim}
realize(LocalConnectivity, WiredEthernet)
\end{verbatim}

B1:

\begin{verbatim}
realize(LocalConnectivity, WiFi)
\end{verbatim}

Projection F ignores realization technology under its declared scope and
can rebuild with the same flow meaning.

Projection A includes realization/assurance context and changes local
output.

PASS: different projection deltas are expected because coverage differs.

\subsection{13. M2 replay - responsibility
change}\label{13-m2-replay---responsibility-change}

When transport becomes project-owned:

\begin{itemize}
\tightlist
\item
  negative responsibility is replaced by affirmative responsibility;
\item
  external service consumption may retire;
\item
  Projection F\textquotesingle s external-service exposure rule ceases
  to apply;
\item
  Projection A can emit a different authority-placement item.
\end{itemize}

PASS: local categories evolve through rebuild from BA, not through BA
mutation from the projections.

\subsection{14. M3 replay - remote to local
recognition}\label{14-m3-replay---remote-to-local-recognition}

The remote transfer is retired by BA.

An exhaustive current-flow projection must omit the old exchange in B1
because it is no longer eligible current BA meaning.

A method view retaining the B0 flow merely because its local item
existed fails rebuild semantics.

PASS: no cross-baseline projection item identity is required.

\subsection{15. M4 replay - stale/unresolved
representation}\label{15-m4-replay---staleunresolved-representation}

A current projection excludes stale capture transfer under its
qualification policy.

A review projection may include:

\begin{verbatim}
stale transfer
unresolved representation conflict
\end{verbatim}

with explicit qualification.

REJECTED controls:

\begin{itemize}
\tightlist
\item
  silently rewrite transfer payload from sibling Decision;
\item
  show stale transfer as current;
\item
  use diagnostic as new governed project fact.
\end{itemize}

\subsection{16. Order/WMS replay}\label{16-orderwms-replay}

Externalizing inventory authority toward a WMS produces different local
effects:

\begin{verbatim}
flow/integration projection
  new service/interchange focus

assurance/responsibility projection
  moved authority/contract dependency focus
\end{verbatim}

Unrelated payment/fulfillment BA meaning is not invalidated merely
because one projection changes strongly.

PASS: method deltas remain localized by BA trace and coverage.

\subsection{17. Provider-normalization
replay}\label{17-provider-normalization-replay}

A method that needs \texttt{PaymentAuthorizationResult} must consume the
accepted BA normalized result.

The method may then derive a local branch or assurance category under
its own interpretation rule.

REJECTED:

\begin{verbatim}
raw provider state
  -> private method mapping
  -> projected value presented as shared governed result
\end{verbatim}

PASS: BA3 remains the authority-preserving normalization layer.

\subsection{18. Cross-baseline and cross-projection
comparison}\label{18-cross-baseline-and-cross-projection-comparison}

T2 tries to force a new shared identity for projection items across
methods and baselines.

The requirement does not survive.

Same-baseline comparison uses role-bound BA trace. Cross-baseline
comparison first resolves BA3 continuity, then rebuilds each projection
against the target BA baseline.

Projection-local continuity may still be useful to a UI/tool, but
current DDTA evidence does not require it as a Base Analysis projection
semantic responsibility.

\subsection{19. Attempts to force earlier
reopen}\label{19-attempts-to-force-earlier-reopen}

\subsubsection{BA1}\label{ba1}

No method-local node/category requires new shared project-semantic
identity with subtype-specific invariants.

\texttt{BA1\ reopen\ =\ NOT\ TRIGGERED}.

\subsubsection{BA2}\label{ba2}

No method-owned category exposes missing methodology-neutral project
meaning requiring a new operator/role.

\texttt{BA2\ reopen\ =\ NOT\ TRIGGERED}.

\subsubsection{BA3}\label{ba3}

Source provenance, review/freshness and BA continuity remain sufficient
to qualify projection input and rebuild after change.

\texttt{BA3\ reopen\ =\ NOT\ TRIGGERED}.

\subsection{20. T2 dispositions}\label{20-t2-dispositions}

\begin{verbatim}
Two incompatible method taxonomies                       PASS
Shared projection ontology                               REJECTED
Method interpretation promoted into BA                   REJECTED
Descriptor + trace alone for multiple local rules        REJECTED
Interpretation rule reference                            REQUIRED
Role-bound BA trace                                      RETAINED
Selective/exhaustive coverage collapse                   REJECTED
Coverage mode                                            REQUIRED
Exhaustive omission -> project negation                  REJECTED
Exhaustive omission -> projection defect                 REQUIRED
Qualification policy                                     REQUIRED
Diagnostic/stale review projection                       PASS
Non-current BA presented as current shared truth         REJECTED
Cross-projection taxonomy equality                       NOT REQUIRED
Cross-projection comparison through BA trace             PASS
New shared projection identity                           NOT FORCED
Projection item BA-like lifecycle                        NOT FORCED
Universal projection DSL                                 REJECTED
New BAE family                                           NOT FORCED
New BA2 operator                                         NOT FORCED
BA1 / BA2 / BA3 reopen                                   NOT TRIGGERED
BA4-T1                                                   NOT REOPENED
BA4                                                      STARTED / NOT CLOSED
\end{verbatim}

\subsection{21. Provisional conclusion}\label{21-provisional-conclusion}

BA4-T2 does not falsify the T1 projection boundary, but it does falsify
two under-specified parts of the lower bound:

\begin{enumerate}
\def\labelenumi{\arabic{enumi}.}
\tightlist
\item
  method-owned interpretation needs item-to-rule accountability when
  several local rules can use overlapping BA bases; and
\item
  coverage needs to distinguish selective views from views that promise
  exhaustive representation within an inspectable declared BA scope.
\end{enumerate}

The accepted refinement remains projection-local and does not expand BA
ontology.

\subsection{22. Smallest remaining BA4
question}\label{22-smallest-remaining-ba4-question}

The next step should be an integrated closure review, not another method
taxonomy exercise.

It must try to remove the T2 refinements, regress the combined T1/T2
contract across both corpora and determine whether BA4 can close for
current thesis scope.

\subsection{23. Next authorized
microstep}\label{23-next-authorized-microstep}

Only after this T2 package is reviewed, committed, pushed and remotely
verified:

\begin{quote}
\textbf{BA4-T3 - projection boundary, interpretation/coverage and
cross-projection closure review.}
\end{quote}

Do not begin BA5 before BA4 is explicitly closed.

\end{document}
